%!TEX root = ../../main.tex
%% ===============================================================
%% 2.8 데이터 누출, 공변량 이동, 선택 편향
%% 파일: chapters/02_background/08_leakage_shift_bias.tex
%% 상위: chapters/02_background/chapter.tex
%% 정의 라벨: sec:bg-leakage
%% 참조 라벨: sec:problem-contract
%% 포함 객체: -
%% 인용: berkson1946limitations, kaufman2012leakage, quinonero2009dataset, shimodaira2000improving
%% 상태: 초안 (docs/OUTLINE.md 참고)
%% ===============================================================

\section{데이터 누출, 공변량 이동, 선택 편향}
\label{sec:bg-leakage}

\begin{definition}[데이터 누출]
학습 시점에 정당하게 사용할 수 없는 정보(예: 예측 대상 사건 이후의 상태)가 입력에 포함되어 모델 성능이 과대평가되는 현상 \cite{kaufman2012leakage}.
\end{definition}

\begin{definition}[공변량 이동]
입력 분포 $P(X)$는 학습과 시험 사이에서 달라지지만 조건부 분포 $P(Y \mid X)$는 유지되는 분포 이동 \cite{shimodaira2000improving,quinonero2009dataset}. 본 논문에서는 패치가 이동의 원인이며, LightGBM에 최근성 가중치를 부여하여 완화한다.
\end{definition}

\begin{definition}[Berkson의 역설]
두 변수가 모집단에서는 독립이더라도, 둘 중 하나 이상에 조건화된 부분 집단만 관측하면 가짜 상관이 생기는 현상 \cite{berkson1946limitations}. 교전 검출 규칙이 예측 지평의 관측 가능한 신호(예: 킬)를 조건으로 하면 $P(Y\mid X, \text{signal}) \ne P(Y \mid X)$가 되어 라벨 분포가 왜곡된다.
\end{definition}

본 논문은 교전의 \emph{존재}를 킬로 조건화하되(킬 조건부 정의), 라벨 창 안의 \emph{결과}에는 조건화하지 않으며, 관측 창의 모든 입력이 온셋 이전 정보만을 사용하도록 계약을 명시한다(제 \ref{sec:problem-contract} 절).
